\section{Companion lab: draw a square using coordinate feedback}
\label{revision-lab:marsbot}

The display counter issues a store for each requested glyph. A moving device
needs a longer-lived command: start it, observe its progress, and stop it when a
condition holds. The command does not finish merely because its store has
completed. This lab uses that distinction to draw a square with MarsBot.

\begin{samepage}
Use \texttt{marsbot-square.asm} in \texttt{revision/labs/files}, the revised work
area for original project \texttt{006-marsbot-square}. Implement the program
before its exit code, replacing the failing \texttt{break}. It receives a word
\texttt{side} from one through 63. Starting at $(0,0)$, visit
\[
 (side,0),\quad(side,side),\quad(0,side),\quad(0,0)
\]
in that order with the pen down, then finish stopped at the origin, facing north,
with the pen up. The trail must contain exactly those four sides. A saved list of
intended corners does not establish that the robot actually traveled between them.\par
\end{samepage}

The prerequisites are the counter's conditional branches, loop state and
memory-mapped stores. A complete one-leg experiment below separates the movement
protocol from your four-leg implementation. You do not need function calls,
arrays or a command table to complete the square. A later table-driven version
is an optional way to revisit it after learning those subjects.

\subsection*{Distinguish command registers from feedback}

The source header connects MarsBot and selects 64 completed simulator instructions
per movement tick. Use MIPS64 Release 6 and the Default memory layout. In that
layout its register block begins at \texttt{0x00000000FFFF8000}. The explicit
leading zeros matter: this is a positive address in the simulated memory map,
not a negative signed displacement from zero. The supplied experiment loads it
with \texttt{dli}. Keep a full address in the base register when adding offsets.

\begin{table}[htbp]
\centering
\small
\begin{tabular}{rlll}
\toprule
Offset & Register & Program access & Meaning\\
\midrule
\texttt{0x10} & heading & word store & Degrees clockwise from north\\
\texttt{0x20} & pen & word store & 1 down, 0 up\\
\texttt{0x30} & $x$ feedback & word load & Current horizontal coordinate\\
\texttt{0x40} & $y$ feedback & word load & Current vertical coordinate\\
\texttt{0x50} & movement & word store & Nonzero moving, zero stopped\\
\bottomrule
\end{tabular}
\caption{MarsBot offsets from the Default-layout base. Feedback reports the device's position; it is not an input for assigning a position. Use a fresh session for each independent route.}
\label{tab:revision-marsbot-ports}
\end{table}


A command is a store to heading, pen or movement control. Feedback is a load from
one of the coordinate registers. Writing a coordinate does not move the robot
there; the device republishes its actual position. This assignment forbids those
stores, including writes later restored. Read feedback to determine where the
robot has reached.

Heading zero points north, 90 east, 180 south and 270 west. Coordinates increase
rightward and downward, so moving south increases $y$. The board has positions
zero through 511 on each axis and clamps motion at its edges. At the initial
upper-left corner, northward motion cannot produce an outward first side.
Figure~\ref{fig:revision-marsbot-route} shows the required route in these coordinates.

\begin{figure}[htbp]
\centering
\begin{tikzpicture}[x=32mm,y=-32mm,alt={A clockwise square in screen coordinates. Origin zero comma zero is upper left. East leads to side comma zero at upper right, south to side comma side at lower right, west to zero comma side at lower left, and north returns to the origin. The x axis points right and y points down.}]
\draw[->,thin] (-.25,-.55)--(1.45,-.55) node[right] {$x$};
\draw[->,thin] (-.9,-.25)--(-.9,1.4) node[below] {$y$};
\draw[->,very thick] (0,0)--(.9,0);
\draw[->,very thick] (1,0)--(1,.9);
\draw[->,very thick] (1,1)--(.1,1);
\draw[->,very thick] (0,1)--(0,.1);
\node[anchor=south east,font=\small] at (0,0) {$(0,0)$};
\node[anchor=south west,font=\small] at (1,0) {$(side,0)$};
\node[anchor=north west,font=\small] at (1,1) {$(side,side)$};
\node[anchor=north east,font=\small] at (0,1) {$(0,side)$};
\node[font=\small,fill=white] at (.5,0) {East 90};
\node[font=\small,fill=white,rotate=90] at (1,.5) {South 180};
\node[font=\small,fill=white] at (.5,1) {West 270};
\node[font=\small,fill=white,rotate=90] at (0,.5) {North 0};
\end{tikzpicture}
\caption{The required route uses screen coordinates. Heading zero points toward decreasing $y$; it is the returning direction, not the first outward leg. Arrows describe the route, not host-time animation.}
\label{fig:revision-marsbot-route}
\end{figure}


Begin each leg while stopped. Set its heading, then store a nonzero movement
command. Stop by storing zero to movement control. The pen state is independent:
starting movement with the pen up changes position without drawing a stroke;
stopping with the pen down does not lift it. These distinctions are part of the
requested final state as well as the route.

\subsection*{Use a reached coordinate as the stopping condition}

The complete \texttt{marsbot-east.asm} experiment draws only the first side.
It loads \texttt{side}, lowers the pen, selects heading 90 and starts movement.
The following loop observes $x$ until it equals the requested endpoint:

\noindent\begin{minipage}{\linewidth}
\MARSCodeLabel{Read the endpoint, then stop the movement command}
\begin{minted}{mars-idea}
east_wait:
    lw $t2, 0x30($s0)
    bnec $t2, $s1, east_wait
    nop
east_stop:
    sw $zero, 0x50($s0)
\end{minted}
\end{minipage}

Here \texttt{\$s0} holds the device base and \texttt{\$s1} the positive side
length. A branch mismatch returns to the load, obtaining fresh feedback. Reusing
the same \texttt{\$t2} value without another load would keep testing an old
observation even while the device moves. The \texttt{nop} keeps the compact
conditional branch's forbidden slot safe; the stop store follows it when the
comparison matches.

\begin{samepage}
During this first leg, the target and heading stay fixed, $y=0$, and $x$ advances
toward the target. The loop must stop movement before another tick can carry it
past the endpoint. Repeatedly reading a reached coordinate is not a stop command.
A picture with the robot at the corner while its status still says Moving is an
unfinished transition.\par
\end{samepage}

This polling argument depends on both the device and the loop. For this
assignment, the rate is an integer from eight through 128 completed instructions
per tick. The short load/branch loop and immediate stop fit the response window.
Do not add printing, long calculations or a second busy-wait between the load
and stop. A different device rate or longer control loop needs a new latency
argument; equality alone cannot recover an endpoint already skipped.

The full \texttt{marsbot-east.asm} program then lifts the pen and exits at
$(side,0)$, still facing east. That is correct for the experiment, but deliberately
incomplete for the square assignment. Its public
\texttt{revision/labs/checks/marsbot-east.kts} check verifies the straight stroke,
its stopped endpoint, protected data and final pen-up state at several lengths
and rates.

\subsection*{Trace the command before and after it takes effect}

Use a fresh Debug launch of the east experiment. At \texttt{east\_wait}, read the
feedback and predict whether the branch will return to the load. Run to
\texttt{east\_stop}: the branch has matched, but the next instruction has not yet
stopped movement. The endpoint is reached and the pen is down; movement remains
on. Step once to \texttt{east\_stopped}. Position is unchanged, but movement is
now off. The pending instruction lifts the pen. At \texttt{leg\_finished}, that
pen command has completed too.

Figure~\ref{fig:revision-marsbot-stop} compares these two movement states for
side 48, used to make the trail easier to see. The supplied source still defaults
to side 16; change only that input to reproduce the illustrated case.

\begin{figure}[htbp]
\centering
\begin{tikzpicture}[alt={Two actual native MarsBot observations for side48 at east_stop and after one native Debug step at east_stopped. Both show position48 comma zero, heading90 degrees and pen down. A native publication box marks Moving in the first and Stopped in the second. Enlarged trail strips below show the same horizontal stroke with its east-facing robot marker. Endpoint labels are native LaTeX annotations.}]
\node[anchor=north west,inner sep=0,font=\small\bfseries] (atitle) at (0,0) {1. Before the stop store};
\node[anchor=north west,inner sep=0] (a) at ([yshift=-2mm]atitle.south west) {\MARSRevisionCrop{revision/captures/marsbot/bot-east-moving.png}{0.43\linewidth}{0.640625000}{0.655555556}{0.195312500}{0.277777778}};
\node[anchor=north west,inner sep=0,font=\small\bfseries] (btitle) at ([xshift=8mm]atitle.north west -| a.north east) {2. After one Debug step};
\node[anchor=north west,inner sep=0] (b) at (btitle.south west |- a.north west) {\MARSRevisionCrop{revision/captures/marsbot/bot-east-stopped.png}{0.43\linewidth}{0.640625000}{0.655555556}{0.195312500}{0.277777778}};
\foreach \panel in {a,b} {
\begin{scope}[shift={(\panel.south west)},x={($(\panel.south east)-(\panel.south west)$)},y={($(\panel.north west)-(\panel.south west)$)}]
\draw[black,line width=.7pt] (.02,.08) rectangle (.70,.42);
\end{scope}}
\node[anchor=north west,inner sep=0] (atr) at ([yshift=-2mm]a.south west) {\MARSRevisionCrop{revision/captures/marsbot/bot-east-moving.png}{0.38\linewidth}{0.646875000}{0.536666667}{0.310156250}{0.441111111}};
\node[anchor=north west,inner sep=0] (btr) at ([yshift=-2mm]b.south west) {\MARSRevisionCrop{revision/captures/marsbot/bot-east-stopped.png}{0.38\linewidth}{0.646875000}{0.536666667}{0.310156250}{0.441111111}};
\foreach \panel in {atr,btr} {
\begin{scope}[shift={(\panel.south west)},x={($(\panel.south east)-(\panel.south west)$)},y={($(\panel.north west)-(\panel.south west)$)}]
\node[anchor=north,font=\small] at (.145,-.02) {$(0,0)$};
\node[anchor=north,font=\small] at (.803,-.02) {$(48,0)$};
\end{scope}}
\end{tikzpicture}
\caption{The endpoint is reached before the stop command executes. One native step changes the movement state while preserving position, pen and completed stroke. Gray lines belong to the board border. The strips enlarge the real trail; boxes and endpoint labels are publication annotations. The 48-pixel example is within the same one-through-63 input contract.}
\label{fig:revision-marsbot-stop}
\end{figure}


Leave the program paused briefly, then inspect the same coordinates. MarsBot
advances on completed simulator instructions, not on a host-time animation timer.
Waiting at a breakpoint executes no instructions and produces no further motion.
Stepping does advance the simulated instruction clock. A slow visual repaint is
not a separate instruction or an extra device tick.

Now change \texttt{instructionsPerTick} in the source header from 64 to 17 and
restart. The device needs fewer completed instructions to move one pixel, so the
program reaches the same coordinate after less polling work. The endpoint and
stroke should remain the same. Keep the illustrated 64 setting when comparing
an intermediate Debug state with a figure; changing the rate can change the
instruction at which a particular coordinate first becomes visible.

A fixed iteration count makes a different promise. It can wait approximately
the right number of instructions at one rate and cover a different distance at
another. That is why the assignment checks actual coordinate feedback and trail
behavior at changed rates. The 100,000-instruction ceiling catches nontermination;
it is not an optimization score or a requirement to match a reference count.

\subsection*{Compose four stopped corners}

Implement the remaining legs in the independent square work area. Set the east
heading first and poll $x$ for \texttt{side}. After stopping, set south and poll
$y$ for \texttt{side}. Next move west and poll $x$ for zero, then move north and
poll $y$ for zero. The other coordinate remains fixed along each axis-aligned
leg. A zero-target returning leg still needs a start command and fresh feedback;
it cannot reuse the nonzero outward target.

\begin{table}[htbp]
\centering
\small
\begin{tabular}{lllll}
\toprule
Completed leg & Position & Heading & Movement & Pen\\
\midrule
East & $(side,0)$ & 90 & stopped & down\\
South & $(side,side)$ & 180 & stopped & down\\
West & $(0,side)$ & 270 & stopped & down\\
North & $(0,0)$ & 0 & stopped & down\\
Final pen command & $(0,0)$ & 0 & stopped & up\\
\bottomrule
\end{tabular}
\caption{A stopped-corner trace. Reaching the final coordinate completes the geometry; the final pen-up command completes the requested device state.}
\label{tab:revision-marsbot-corners}
\end{table}


For this construction, stop before changing heading. This gives each corner a boundary at which position,
heading and trail have a clear interpretation. Keeping the pen down through a
stopped corner preserves the continuous square; lift it after the final stop.
Heading north on the last leg already establishes the required final heading.
There is no need to draw an extra turn or move beyond the origin. The assessed
contract is the complete trail and final state; it does not prescribe a particular
intermediate command count.

Use \texttt{side} throughout. With length one, each leg is only one device tick,
which makes command/stop latency especially visible. With length 63, the robot
travels farther but the same four transitions and tests apply. None of the allowed
squares reaches the opposite boundary of the 512-position board, so a boundary
clamp cannot substitute for the outward endpoint test.

If you use four explicit legs, the target, feedback offset and heading should be
easy to identify in each. After learning array traversal, you can represent those
three choices in a command table and drive one common polling loop. Both forms
can satisfy the same route contract. The grader does not prescribe private labels,
scratch registers or an exact command count.

\subsection*{Diagnose a route, not only a final position}

Begin with the default length 16. Record the state at each stopped corner and
compare the accumulated trail with the table. Repeat with length one and another
nondefault input. Then run the completed square at a changed instruction rate.
The final position must agree, but so must every straight side and the final
heading, movement and pen states.

In a disposable copy, leave the pen up for the first side. The robot can still
visit every corner and return to the origin, but the first stroke is missing.
In another copy, replace the input load with a constant 16. The default drawing
looks correct; a length-one check exposes that the input is ignored. Restore each
fault before introducing another.

For a feedback fault, use the $x$ register in the southward loop. The fixed $x$
value can make that loop stop immediately instead of drawing the second side.
For a control fault, omit the final stop before lifting the pen. The robot can
remain at the origin because the board blocks further northward motion, while
its movement state is still on. Inspect that state before exit. Neither fault is
diagnosed by looking only at the final origin. Keep the first divergent command,
stroke or device-state field in your report.

The trail must contain exactly the four requested straight sides, with no pen-up
gaps or extra strokes. The program must not store to \texttt{side},
\texttt{side\_guard}, or either coordinate register, even temporarily, and must
produce no console output. Device feedback reads and control stores have different
roles; preserved RAM words do not replace a real device observation.

\begin{samepage}
Run the original public check on your completed square:
\begin{verbatim}
java -jar /path/to/cli.jar grade \
  packages/006-marsbot-square/checks/public.kts \
  /path/to/your/marsbot-square.asm
\end{verbatim}
\par\end{samepage}

The cases change length and movement rate, inspect the actual device trail and
final state, and reject protected stores. Keep your corner trace and one polling
invariant with the passing result. Explain separately why the robot reached each
corner, why it stopped there, and why its path contains the required ink.
